\section{Introduction}
\label{sec:introduction}

According to the United Nations Environment Programme, households
account for the largest share of total food waste worldwide, with
recent figures attributing more than 60\% of global food waste to the
household stage of the supply chain~\cite{unep2024foodwaste}. Beyond
its environmental cost, this loss directly inflates household
expenditure on groceries, with FAO estimating the economic value of
wasted food in the range of one trillion US dollars per year. Food
that is purchased, refrigerated, then thrown away after spoiling is
particularly costly because it has consumed packaging, transport, and
cooling resources up to the point of disposal.

The most common reasons for in-home food waste are well documented:
users (i) forget what is already in the refrigerator,
(ii) underestimate or misremember expiration dates, and
(iii) only notice spoilage once a product has visually deteriorated.
Manual logging applications exist, but their adoption is limited
because they require disciplined data entry and do not give visibility
into items already in storage.

Several smart-refrigerator and inventory-monitoring approaches have
been proposed in the literature to mitigate these issues.
Camera-based product recognition systems can in principle identify
items automatically but require trained classification models, generate
substantial computational load on the edge device, and are sensitive to
lighting and partial occlusion~\cite{ericyang2022smartfridge}.
Radio-frequency identification (RFID) approaches eliminate the manual
data-entry step but assume every monitored product carries a tag and
require a reader installed in the storage area, which raises both unit
cost and integration complexity~\cite{abhiram2025iotmonitor}.
Barcode-based systems are inexpensive but inherently manual and assume
unobstructed line of sight to a printed
label~\cite{pelobello2026chickout}. Many recent prototypes focus on a
single sensing dimension --- temperature, weight, or visual recognition
--- and rarely fuse multiple modalities or expose a real-time mobile
interface to the user~\cite{sarawathi2025spoilage,sarkar2024proactive}.

This work proposes the \emph{Smart Inventory and Automated Food
Tracking Module}, a low-cost prototype that addresses the limitations
above by combining four complementary sensing modalities ---
weight, temperature, humidity, gas, and visual --- with QR-based
product registration and a real-time mobile interface in a single
cabinet enclosure. An Arduino Uno serves as the primary microcontroller
for the load-cell, DHT11, and MQ135 subsystems; an AI Thinker ESP32-CAM
hosts a small HTTP server providing a live MJPEG feed and JPEG capture
for QR scanning; a Python USB-serial bridge exposes the sensor stream
as a local HTTP endpoint; and a cross-platform Flutter application
visualizes the entire system, registers products, and surfaces alerts.

\subsection*{Contributions}

The principal contributions of this work are:

\begin{itemize}
  \item A self-contained, low-cost cabinet prototype that simultaneously
        provides shelf-weight measurement, temperature and humidity
        monitoring, indicative gas-trend monitoring, QR-based product
        registration, and LED-assisted illumination, all coordinated
        through a single Arduino Uno node and an auxiliary ESP32-CAM
        node.
  \item A Python USB-serial-to-HTTP bridge that exposes the Arduino
        Uno's sensor output as a CORS-enabled local web endpoint,
        eliminating the need for a Wi-Fi-capable sensor
        microcontroller and enabling browser- and mobile-side
        consumption without firmware changes.
  \item A cross-platform Flutter mobile dashboard (Android and web
        targets) with five screens --- Dashboard, Camera, Products,
        Alerts, and Settings --- that integrates the bridge feed and
        the camera's MJPEG stream, performs in-frame multi-QR
        decoding, tracks products with expiration metadata, and emits
        rule-based alerts.
  \item A transparent, pixel-level RGB classifier for banana surface
        spotting that runs on commodity hardware and reports a single
        interpretable percentage --- an honest baseline against which
        machine-learning-based approaches can be compared.
  \item A complete, openly archived implementation (firmware,
        bridge, application, and documentation), preserved in the
        Software Heritage Archive with a persistent identifier so that
        the prototype is reproducible by independent groups.
\end{itemize}

The remainder of this paper is organized as follows.
Section~\ref{sec:relatedwork} reviews related work on smart-refrigerator
and inventory systems and positions the proposed module against the
state of the art. Section~\ref{sec:systemdesign} describes the
high-level system architecture and the rationale for each subsystem
choice. Section~\ref{sec:implementation} details the hardware
configuration, firmware, USB-serial bridge, and mobile application.
Section~\ref{sec:algorithms} presents the QR-handling, ripeness, and
risk-score algorithms. Section~\ref{sec:results} reports qualitative
results obtained during demonstration sessions.
Section~\ref{sec:discussion} discusses the implications and the
limitations that should be addressed in future work, and
Section~\ref{sec:conclusion} concludes.
